\input{../tex-inputs/latex-korrekturansicht-vorspann}

\section[Arthur Schnitzler an Gustav Schwarzkopf, 26. 10. 1922]{L04363 Arthur Schnitzler an Gustav Schwarzkopf, 26. 10. 1922}
\nopagebreak\mylabel{L04363v}
\rehead{ }\normalsize\beginnumbering\briefempfaengerindex{Schwarzkopf, Gustav@\textsc{Schwarzkopf, Gustav}!zzzSchnitzler, Arthur@\emph{von Arthur Schnitzler}!1922-10-261@{26. 10. 1922}|(be}
\toendnotes[C]{\smallbreak\pagebreak[2]}
\correspDesc{Versand  durch Arthur Schnitzler am 26. 10. 1922 in Brünn
\newline{}Erhalt  durch Gustav Schwarzkopf im Zeitraum [27. 10. 1922 – 31. 10. 1922?] in Wien}\toendnotes[C]{\smallbreak}
\Standort{DLA, A:Schnitzler, HS.1985.1.1897.}
\physDesc{Bildpostkarte, 141 Zeichen
\newline{}Handschrift: Bleistift, deutsche Kurrent
\newline{}Versand: Stempel: »\nobreak{}\oindex{Brünn@\textbf{Brünn}|pwk}Brno 2, 26. X 2\textcolor{gray}{2}, 9\nobreak{}«.  }\pstart{}{\pb}Hr Guſtav Schwarzkopf\pend{}\pstart{}\textcolor{pink}{Wien I}\oindex{I., Innere Stadt@\textbf{I., Innere Stadt}, \emph{Verwaltungsgebiet}|pw}{}\ledrightnote{\textcolor{pink}{I., Innere Stadt}}\pend{}\pstart{}\textcolor{pink}{Tiefer Graben 17}\oindex{Wien@\textbf{Wien}!I., Innere Stadt@\textbf{I., Innere Stadt}!Tiefer Graben 17@\textbf{Tiefer Graben 17}, \emph{Wohngebäude}|pw}{}\ledrightnote{\textcolor{pink}{Tiefer Graben 17}}\pend{}{\bigskip}
\pstart
           \noindent{}\centering{}{\pb}\textcolor{pink}{\textcolor{gray}{\textbf{Brünn. Altbrünn}}}\oindex{Staré Brno@\textbf{Staré Brno}, \emph{Teil eines besiedelten Ortes}|pw}{}\ledrightnote{\textcolor{pink}{Staré Brno}}\pend
           \vspace{1em}
\pstart
           \noindent{}{\pb}\textsc{Nach guter alter Sitte,}{ }\textsc{lieber Guſtav,}{ }\textsc{auch von hier aus die herzlichſten Grüße!}\pend
           
\pstart
           \textsc{Ihr}{\\[\baselineskip]}\spacefill\mbox{A.}\pend
           \leftskip=0em{}
\pstart
           \textcolor{pink}{Brünn}\oindex{Brünn@\textbf{Brünn}|pw}{}\ledrightnote{\textcolor{pink}{Brünn}}{ }26/X 22\pend
           \selectlanguage{ngerman}\endnumbering\briefempfaengerindex{Schwarzkopf, Gustav@\textsc{Schwarzkopf, Gustav}!zzzSchnitzler, Arthur@\emph{von Arthur Schnitzler}!1922-10-261@{26. 10. 1922}|)be}\mylabel{L04363h}
\begin{anhang}
\end{anhang}\newcommand{\dateiname}{L04363}\newcommand{\titel}{Arthur Schnitzler an Gustav Schwarzkopf, 26. 10. 1922}\newcommand{\editorInnen}{Herausgegeben von Jahnke, SelmaMüller, Martin Anton}\input{../tex-inputs/latex-korrekturansicht-abspann}
